% Vietnamese translation for OI Public Library
% Source-PDF: IOI中国国家候选队论文/2002/何江舟.pdf
\documentclass[11pt]{article}
\usepackage{oipl-vn}

\title{Dùng khử Gauss để giải hệ phương trình tuyến tính}
\author{He Jiangzhou}
\date{}

\begin{document}
\maketitle

\origtitle{Using Gaussian elimination to solve systems of linear equations}
\sourcepath{IOI China candidate team papers / 2002 / He Jiangzhou.pdf}

\tableofcontents

\section{Bài toán mở đầu: hệ thống xếp hạng GPA}

Bài toán GPA Ranking System của CTSC 2001 là một ví dụ tự nhiên dẫn đến hệ
phương trình tuyến tính.

Các trường đại học thường dùng GPA để đánh giá kết quả học tập của sinh viên.
Cách xếp hạng truyền thống là tính điểm trung bình của từng sinh viên rồi xếp
hạng theo điểm trung bình đó. Tuy nhiên, với các môn học khác nhau, điểm trung
bình của nhóm sinh viên chọn môn đó chịu ảnh hưởng bởi độ khó của môn và nhiều
yếu tố khác, nên cách xếp hạng trực tiếp theo trung bình thô chưa thật hợp lý.

Ta muốn cải tiến hệ thống xếp hạng như sau. Với mỗi môn học thứ \(i\), cộng
vào điểm của mọi sinh viên học môn đó một giá trị hiệu chỉnh riêng \(d_i\).
Điểm sau hiệu chỉnh không nhất thiết còn nằm trong thang \(100\). Các giá trị
\(d_i\) phải được chọn sao cho sau khi hiệu chỉnh, điểm trung bình của riêng
môn đó bằng điểm trung bình tất cả các môn của tất cả sinh viên đã chọn môn
đó. Điều kiện này phải đúng với mọi môn học. Nhiệm vụ là dựa vào điểm của một
khối sinh viên trong một năm học, thực hiện hiệu chỉnh như trên rồi đưa ra xếp
hạng.

Ký hiệu:
\[
\begin{array}{ll}
A_i & \text{tập sinh viên chọn môn thứ } i,\\
B_j & \text{tập môn học mà sinh viên thứ } j \text{ chọn},\\
G_{i,j} & \text{điểm của sinh viên thứ } j \text{ ở môn thứ } i,\\
d_i & \text{giá trị hiệu chỉnh của môn thứ } i.
\end{array}
\]

Với môn thứ \(p\), điều kiện cân bằng trung bình được viết thành
\[
  \frac{1}{|A_p|}\sum_{j\in A_p}\bigl(G_{p,j}+d_p\bigr)
  =
  \frac{1}{\sum_{j\in A_p}|B_j|}
  \sum_{j\in A_p}\sum_{i\in B_j}\bigl(G_{i,j}+d_i\bigr).
\]
Đây là một phương trình tuyến tính theo các ẩn
\(d_1,d_2,\ldots,d_n\). Viết phương trình như vậy cho từng môn, ta thu được
một hệ gồm \(n\) phương trình tuyến tính.

\section{Dạng tổng quát của hệ phương trình tuyến tính}

Một hệ tuyến tính \(n\) ẩn có dạng
\[
\begin{aligned}
  a_{1,1}x_1+a_{1,2}x_2+\cdots+a_{1,n}x_n&=b_1,\\
  a_{2,1}x_1+a_{2,2}x_2+\cdots+a_{2,n}x_n&=b_2,\\
  &\vdots\\
  a_{n,1}x_1+a_{n,2}x_2+\cdots+a_{n,n}x_n&=b_n.
\end{aligned}
\]
Ta thường biểu diễn hệ này bằng ma trận mở rộng
\[
\left[
\begin{array}{cccc|c}
a_{1,1} & a_{1,2} & \cdots & a_{1,n} & b_1\\
a_{2,1} & a_{2,2} & \cdots & a_{2,n} & b_2\\
\vdots  & \vdots  &        & \vdots  & \vdots\\
a_{n,1} & a_{n,2} & \cdots & a_{n,n} & b_n
\end{array}
\right].
\]
Khử Gauss biến ma trận mở rộng này về dạng tam giác trên, rồi dùng hồi thế để
tính nghiệm.

\section{Một ví dụ}

Xét ma trận mở rộng
\[
\left[
\begin{array}{rrr|r}
2 & -1 & 3 & 1\\
4 & 2 & 5 & 4\\
1 & 2 & 0 & 7
\end{array}
\right].
\]

Ở bước đầu, dùng phần tử chính \(2\) của hàng thứ nhất để khử các phần tử bên
dưới nó. Hàng thứ hai trừ \(2\) lần hàng thứ nhất; hàng thứ ba trừ \(0.5\) lần
hàng thứ nhất. Ta được
\[
\left[
\begin{array}{rrr|r}
2 & -1 & 3 & 1\\
0 & 4 & -1 & 2\\
0 & 2.5 & -1.5 & 6.5
\end{array}
\right].
\]
Tiếp theo dùng phần tử chính \(4\) của hàng thứ hai. Hàng thứ ba trừ
\(2.5/4\) lần hàng thứ hai:
\[
\left[
\begin{array}{rrr|r}
2 & -1 & 3 & 1\\
0 & 4 & -1 & 2\\
0 & 0 & -0.875 & 5.25
\end{array}
\right].
\]

Hồi thế cho kết quả
\[
  x_3=\frac{5.25}{-0.875}=-6,
\]
\[
  x_2=\frac{2-(-1)x_3}{4}=-1,
\]
\[
  x_1=\frac{1-(-1)x_2-3x_3}{2}=9.
\]

\section{Quá trình khử}

Dùng chỉ số trên \((k)\) để biểu diễn trạng thái trước lần khử thứ \(k\). Ban
đầu, trước lần khử thứ nhất, ma trận mở rộng là
\[
\left[
\begin{array}{cccc|c}
a_{1,1}^{(1)} & a_{1,2}^{(1)} & \cdots & a_{1,n}^{(1)} & b_1^{(1)}\\
a_{2,1}^{(1)} & a_{2,2}^{(1)} & \cdots & a_{2,n}^{(1)} & b_2^{(1)}\\
\vdots        & \vdots        &        & \vdots        & \vdots\\
a_{n,1}^{(1)} & a_{n,2}^{(1)} & \cdots & a_{n,n}^{(1)} & b_n^{(1)}
\end{array}
\right].
\]

Ở lần khử thứ nhất, với \(i=2,3,\ldots,n\), hệ số nhân của hàng thứ nhất là
\[
  m_{i,1}=\frac{a_{i,1}^{(1)}}{a_{1,1}^{(1)}}.
\]
Sau khi dùng hàng thứ nhất để khử các phần tử bên dưới cột thứ nhất, với
\(i,j=2,3,\ldots,n\), ta có
\[
  a_{i,j}^{(2)}=a_{i,j}^{(1)}-m_{i,1}a_{1,j}^{(1)},
  \qquad
  b_i^{(2)}=b_i^{(1)}-m_{i,1}b_1^{(1)}.
\]
Ma trận mới có dạng
\[
\left[
\begin{array}{cccc|c}
a_{1,1}^{(1)} & a_{1,2}^{(1)} & \cdots & a_{1,n}^{(1)} & b_1^{(1)}\\
0             & a_{2,2}^{(2)} & \cdots & a_{2,n}^{(2)} & b_2^{(2)}\\
\vdots        & \vdots        &        & \vdots        & \vdots\\
0             & a_{n,2}^{(2)} & \cdots & a_{n,n}^{(2)} & b_n^{(2)}
\end{array}
\right].
\]

Tổng quát hơn, trước lần khử thứ \(k\), ma trận đã có dạng bậc thang
\[
\left[
\begin{array}{cccccc|c}
a_{1,1}^{(1)} & a_{1,2}^{(1)} & \cdots & a_{1,k}^{(1)}
  & \cdots & a_{1,n}^{(1)} & b_1^{(1)}\\
0 & a_{2,2}^{(2)} & \cdots & a_{2,k}^{(2)}
  & \cdots & a_{2,n}^{(2)} & b_2^{(2)}\\
\vdots & & \ddots & \vdots & & \vdots & \vdots\\
0 & 0 & \cdots & a_{k,k}^{(k)}
  & \cdots & a_{k,n}^{(k)} & b_k^{(k)}\\
\vdots & \vdots & & \vdots & & \vdots & \vdots\\
0 & 0 & \cdots & a_{n,k}^{(k)}
  & \cdots & a_{n,n}^{(k)} & b_n^{(k)}
\end{array}
\right].
\]
Phần tử \(a_{k,k}^{(k)}\) là phần tử chính của lần khử thứ \(k\), và hàng thứ
\(k\) là hàng chính. Với \(i=k+1,k+2,\ldots,n\), đặt
\[
  m_{i,k}=\frac{a_{i,k}^{(k)}}{a_{k,k}^{(k)}}.
\]
Khi đó, với \(i,j=k+1,k+2,\ldots,n\),
\[
  a_{i,j}^{(k+1)}
  =a_{i,j}^{(k)}-m_{i,k}a_{k,j}^{(k)},
  \qquad
  b_i^{(k+1)}
  =b_i^{(k)}-m_{i,k}b_k^{(k)}.
\]

\section{Quá trình hồi thế}

Sau khi khử xong, ta thu được ma trận mở rộng tam giác trên
\[
\left[
\begin{array}{ccccc|c}
a_{1,1}^{(1)} & a_{1,2}^{(1)} & \cdots & a_{1,n}^{(1)} & & b_1^{(1)}\\
0 & a_{2,2}^{(2)} & \cdots & a_{2,n}^{(2)} & & b_2^{(2)}\\
\vdots & & \ddots & \vdots & & \vdots\\
0 & 0 & \cdots & a_{n,n}^{(n)} & & b_n^{(n)}
\end{array}
\right].
\]
Từ hàng cuối đi ngược lên, nghiệm được tính bằng
\[
  x_i=
  \frac{
    b_i^{(i)}-\sum_{j=i+1}^{n}a_{i,j}^{(i)}x_j
  }{
    a_{i,i}^{(i)}
  }.
\]

\section{Vì sao cần chọn phần tử chính}

Mô tả ở trên khử theo thứ tự tự nhiên \(x_1,x_2,\ldots,x_n\). Khi đó ở bước
thứ \(k\),
\[
  m_{i,k}=\frac{a_{i,k}^{(k)}}{a_{k,k}^{(k)}}.
\]
Vì vậy phần tử chính của mỗi bước không được bằng \(0\). Nếu trong quá trình
khử tự nhiên xuất hiện \(a_{k,k}^{(k)}=0\), thuật toán không thể tiếp tục.
Ngay cả khi \(|a_{k,k}^{(k)}|\) rất nhỏ, sai số tính toán cũng có thể tăng
mạnh.

Ví dụ, giả sử mọi phép tính đều dùng số thực độ chính xác đơn. Với hệ
\[
\left[
\begin{array}{cc|c}
10^{-10} & 1 & 1\\
1 & 1 & 2
\end{array}
\right],
\]
nếu dùng hàng đầu tiên làm hàng chính, sau bước khử ta nhận được xấp xỉ
\[
\left[
\begin{array}{cc|c}
10^{-10} & 1 & 1\\
0 & -10^{10} & -10^{10}
\end{array}
\right],
\]
và có thể tính ra \(x_1=0,\ x_2=1\). Nghiệm này sai lệch lớn đối với phương
trình thứ hai. Nguyên nhân là hệ số nhân ở bước khử quá lớn; phương trình đầu
đã lấn át phương trình thứ hai, nên thông tin từ phương trình thứ hai gần như
không còn thể hiện trong kết quả.

Nếu đổi thứ tự hai hàng:
\[
\left[
\begin{array}{cc|c}
1 & 1 & 2\\
10^{-10} & 1 & 1
\end{array}
\right],
\]
sau khi khử ta thu được xấp xỉ
\[
\left[
\begin{array}{cc|c}
1 & 1 & 2\\
0 & 1 & 1
\end{array}
\right],
\]
và suy ra \(x_1=1,\ x_2=1\), phù hợp hơn nhiều với hệ ban đầu.

Do đó ở mỗi lần khử, nên chọn phần tử chính có trị tuyệt đối càng lớn càng tốt
để hệ số nhân dùng với hàng chính càng nhỏ càng tốt.

\section{Chọn phần tử chính}

Khi thực hiện lần khử thứ \(k\), xét các phần tử trên cột \(k\) từ hàng \(k\)
trở xuống:
\[
  a_{k,k}^{(k)},a_{k+1,k}^{(k)},\ldots,a_{n,k}^{(k)}.
\]
Chọn phần tử có trị tuyệt đối lớn nhất trong các phần tử này. Nếu nó nằm ở
hàng \(l\), ta đổi hàng \(l\) với hàng \(k\), rồi dùng hàng mới thứ \(k\) làm
hàng chính. Đây là chọn phần tử chính theo cột, hay chọn phần tử chính từng
phần.

\section{Trường hợp vô nghiệm}

Nếu trong quá trình khử xuất hiện một hàng mà mọi hệ số của các ẩn đều bằng
\(0\), nhưng hằng số ở vế phải khác \(0\), tức là một hàng dạng
\[
  \left[\begin{array}{cccc|c}
  0 & 0 & \cdots & 0 & c
  \end{array}\right],
  \qquad c\ne 0,
\]
thì hàng đó biểu diễn mệnh đề \(0=c\). Khi đó hệ phương trình vô nghiệm.

\section{Trường hợp có vô số nghiệm}

Trong quá trình khử cũng có thể xuất hiện một hàng mà mọi hệ số và hằng số đều
bằng \(0\). Điều này tương đương với việc mất đi một phương trình độc lập.
Sau đó số phương trình độc lập nhỏ hơn số ẩn; hệ hoặc vô nghiệm, hoặc có vô
số nghiệm. Với hệ có nghiệm, ta có thể chọn tùy ý các ẩn tự do rồi hồi thế để
lấy một nghiệm cụ thể.

Xét ví dụ
\[
\left[
\begin{array}{rrr|r}
4 & 2 & 3 & 9\\
2 & 1 & 1 & 4\\
2 & 1 & 2 & 5
\end{array}
\right].
\]
Sau khi khử cột thứ nhất, ta được
\[
\left[
\begin{array}{rrr|r}
4 & 2 & 3 & 9\\
0 & 0 & -0.5 & -0.5\\
0 & 0 & 0.5 & 0.5
\end{array}
\right].
\]
Nếu tiếp tục khử cột thứ hai thì cần có phần tử chính \(a_{2,2}\ne 0\), nhưng
cột thứ hai không còn phần tử khác \(0\) nào ở các hàng chưa xử lý. Vì vậy ta
dùng cột thứ ba làm cột chính tiếp theo:
\[
\left[
\begin{array}{rrr|r}
4 & 2 & 3 & 9\\
0 & 0 & -0.5 & -0.5\\
0 & 0 & 0 & 0
\end{array}
\right].
\]
Hàng thứ ba đã mất ý nghĩa ràng buộc. Biến \(x_2\) không được dùng làm biến
chính, nên có thể chọn tùy ý. Chẳng hạn đặt \(x_2=0\), rồi hồi thế được
\(x_3=1\) và \(x_1=1.5\). Đây là một nghiệm khả thi.

Với hệ tuyến tính tổng quát, trước mỗi lần khử ta tìm một cột còn phần tử khác
\(0\) trong các hàng chưa xử lý, dùng phần tử tương ứng làm phần tử chính, rồi
tiếp tục khử. Khi mọi phần tử mới trong các hàng còn lại đều bằng \(0\), ta
chia ẩn thành hai loại. Các ẩn ứng với cột không được chọn làm cột chính là ẩn
tự do và có thể nhận giá trị tùy ý; các ẩn ứng với cột chính được xác định
trong quá trình hồi thế. Nếu sắp xếp lại thứ tự ẩn sao cho các cột chính đứng
trước, thì sau khi có \(k\) cột chính, các ẩn còn lại từ vị trí \(k+1\) đến
\(n\) là ẩn tự do.

\section{Phân tích hiệu năng và sai số}

Độ phức tạp thời gian của khử Gauss là
\[
  O(n^3).
\]
Trong đó, phần khử chiếm \(O(n^3)\), chọn phần tử chính theo cột chiếm
\(O(n^2)\), và hồi thế chiếm \(O(n^2)\).

Độ phức tạp bộ nhớ là
\[
  O(n^2),
\]
chủ yếu dùng để lưu ma trận mở rộng. Nếu dùng chọn phần tử chính toàn phần,
tức là cho phép đổi cả hàng và cột, cần thêm một bảng ghi quan hệ giữa cột và
ẩn, kích thước \(O(n)\).

Về độ chính xác, vì thuật toán dùng phép tính số thực, và mỗi lần khử sau lần
đầu đều sử dụng kết quả được sinh ra từ các lần khử trước, sai số có thể tích
lũy đáng kể. Mỗi bước hồi thế lại dùng cả kết quả khử và các nghiệm đã hồi
thế trước đó, nên phương pháp này thường chỉ cho nghiệm gần đúng. Trong bài
toán cụ thể, cần chọn cách cải tiến phù hợp với yêu cầu về độ chính xác.

\section{Giải chính xác hệ tuyến tính nguyên}

Phần trước xét thuật toán nhận nghiệm gần đúng bằng phép tính số thực. Trong
một số bài toán, ta cần nghiệm chính xác. Ở đây xét trường hợp các hệ số, hằng
số và nghiệm đều là số nguyên.

Ở khử Gauss thông thường, bước khử dùng
\[
  m_{i,k}=\frac{a_{i,k}^{(k)}}{a_{k,k}^{(k)}},
\]
\[
  a_{i,j}^{(k+1)}=a_{i,j}^{(k)}-m_{i,k}a_{k,j}^{(k)},
  \qquad
  b_i^{(k+1)}=b_i^{(k)}-m_{i,k}b_k^{(k)}.
\]
Rõ ràng đây là phép tính trên số thực hoặc phân số.

Vì không thể bảo đảm \(a_{i,k}^{(k)}\) là bội của \(a_{k,k}^{(k)}\), để khử
mà vẫn giữ các phần tử nguyên, ta nhân hai hàng với hai hệ số nguyên thích
hợp. Gọi
\[
  L=\operatorname{lcm}\bigl(|a_{i,k}^{(k)}|,|a_{k,k}^{(k)}|\bigr).
\]
Đặt
\[
  m_{i,k}=\frac{L}{a_{k,k}^{(k)}},
  \qquad
  m'_{i,k}=\frac{L}{a_{i,k}^{(k)}}.
\]
Khi đó hệ số của \(x_k\) trong hàng mới sẽ bị triệt tiêu, và các phần tử còn
lại được cập nhật bằng
\[
  a_{i,j}^{(k+1)}
  =m'_{i,k}a_{i,j}^{(k)}-m_{i,k}a_{k,j}^{(k)},
  \qquad
  b_i^{(k+1)}
  =m'_{i,k}b_i^{(k)}-m_{i,k}b_k^{(k)}.
\]

Khi số phương trình nhiều, các hệ số có thể ngày càng lớn và có nguy cơ vượt
phạm vi biểu diễn. Vì vậy cần thường xuyên giản ước từng hàng: chia mọi phần
tử trong hàng đó cho ước chung lớn nhất của chúng. Dù vậy, cách làm này vẫn
không thể bảo đảm tuyệt đối tránh tràn số. Nó thường phù hợp hơn với các hệ
có hệ số và nghiệm trong phạm vi nhỏ, số ẩn cũng không quá lớn.

\section{Ứng dụng: bánh răng}

Giả sử ta có một bộ đồ chơi gồm nhiều bánh răng với các kích thước khác nhau.
Có nhiều nhất \(20\) loại bánh răng, mỗi loại có vô hạn chiếc, và mỗi bánh
răng có nhiều nhất \(100\) răng. Ta muốn dùng chúng để tạo các bộ truyền động
với nhiều tỉ lệ khác nhau.

Một bộ truyền động gồm số chẵn bánh răng. Các bánh răng được chia thành từng
cặp ăn khớp với nhau, và mỗi cặp được nối với cặp kế tiếp bằng trục. Gọi
\(c_1,c_2,\ldots,c_m\) là số răng của bánh răng thứ nhất trong từng cặp, và
\(d_1,d_2,\ldots,d_m\) là số răng của bánh răng thứ hai trong từng cặp. Khi
đó tỉ lệ truyền động là
\[
  c:d=\prod_{i=1}^{m}\frac{c_i}{d_i}.
\]

Ví dụ, ta có ba loại bánh răng \(6\) răng, \(12\) răng và \(30\) răng, và cần
tạo tỉ lệ \(4:5\). Một phương án khả thi dùng bốn bánh răng, chia thành hai
cặp: cặp thứ nhất là \(12\) răng và \(6\) răng, cặp thứ hai là \(12\) răng và
\(30\) răng. Khi đó
\[
  \frac{12}{6}\cdot\frac{12}{30}=\frac{4}{5}.
\]

Gọi số răng của các loại bánh răng là \(a_1,a_2,\ldots,a_n\). Giả sử loại
\(a_i\) xuất hiện \(e_i\) lần ở vai trò bánh răng phía \(c\), và \(f_i\) lần
ở vai trò bánh răng phía \(d\). Khi đó
\[
  c:d
  =
  \prod_{i=1}^{n}\frac{a_i^{e_i}}{a_i^{f_i}}
  =
  \prod_{i=1}^{n}a_i^{e_i-f_i}.
\]
Một loại bánh răng đồng thời được dùng ở cả hai phía là lãng phí, vì hai lần
xuất hiện đó triệt tiêu nhau trong tích. Đặt
\[
  x_i=e_i-f_i.
\]
Nếu \(x_i>0\), loại bánh răng đó chỉ cần xuất hiện ở phía \(c\); nếu
\(x_i<0\), nó chỉ cần xuất hiện ở phía \(d\). Bài toán chuyển thành tìm các
số nguyên \(x_i\).

Ta phân tích \(c,d\) và các \(a_i\) thành thừa số nguyên tố. Vì \(a_i\le 100\),
các \(a_i\) chỉ có thể chứa các nguyên tố không vượt quá \(100\), tổng cộng
không quá \(25\) loại nguyên tố. Nếu \(c\) hoặc \(d\) chứa một nguyên tố khác
ngoài các nguyên tố này, hiển nhiên không thể tạo được tỉ lệ yêu cầu.

Gọi \(g_{r,i}\) là số mũ của nguyên tố \(r\) trong phân tích của \(a_i\);
gọi \(c_r,d_r\) lần lượt là số mũ của \(r\) trong phân tích của \(c,d\). Khi
đó ta có các quan hệ
\[
  2^{x_1g_{2,1}+x_2g_{2,2}+\cdots+x_ng_{2,n}}=2^{c_2-d_2},
\]
\[
  3^{x_1g_{3,1}+x_2g_{3,2}+\cdots+x_ng_{3,n}}=3^{c_3-d_3},
\]
và tương tự cho các nguyên tố khác. Các quan hệ này tương đương với hệ phương
trình trên số mũ:
\[
\begin{aligned}
  g_{2,1}x_1+g_{2,2}x_2+\cdots+g_{2,n}x_n &= c_2-d_2,\\
  g_{3,1}x_1+g_{3,2}x_2+\cdots+g_{3,n}x_n &= c_3-d_3,\\
  &\vdots\\
  g_{97,1}x_1+g_{97,2}x_2+\cdots+g_{97,n}x_n &= c_{97}-d_{97}.
\end{aligned}
\]
Ngoài ra, vì mỗi cặp bánh răng đóng góp một bánh ở phía \(c\) và một bánh ở
phía \(d\), số bánh ở hai phía phải bằng nhau. Do đó còn có ràng buộc
\[
  x_1+x_2+\cdots+x_n=0.
\]
Như vậy bài toán bánh răng được chuyển hoàn toàn thành bài toán giải hệ
phương trình tuyến tính nguyên.

\section{Kết luận}

Khử Gauss là một thuật toán đơn giản, hiệu quả và có phạm vi áp dụng rộng.
Trong ứng dụng thực tế, ta thường phải phân tích đặc điểm của từng bài toán,
rồi cải tiến thuật toán chuẩn này cho phù hợp với yêu cầu về nghiệm, sai số,
dữ liệu nguyên, hoặc cấu trúc riêng của hệ phương trình.

\end{document}
